\section{Finite free approximations to actual chain homology}

\begin{theorem}[Finite free triples computing each original fibre homology]
\label{thm:kn-derived-finite-free-triples}
\uses{lem:kn-derived-free-replacement-step,lem:kn-derived-free-replacement-lower,lem:kn-derived-free-replacement-upper}
\lean{EnvelopeProperty.Vanishing.SpecializationDerivedFiniteApproximation.exists_finite_free_triple}
\leanok
Let $C_\bullet$ be an actual chain complex in nonnegative degrees over
a Noetherian ring $R$, with flat terms and finite native homology in
every degree.  For each degree $q$, finite free modules and actual
maps between them can be constructed whose native middle homology,
after every $R$-algebra base change, is isomorphic to the native
degree-$q$ homology of the actual scalar-extended original complex.
No comparison or finite-free model is supplied as an input.
\end{theorem}
\begin{proof}
Start with the original degree-zero module and the zero outgoing
term.  Repeatedly apply the constructed finite-free replacement step.
At every stage the new predecessor is flat.  Its incoming map is
the original next differential with zero free component.  Purity of
the actual relation sequence identifies its kernel after every scalar
extension with the kernel of that original next differential.
The lower homology comparison and this kernel identification prove
that the next frontier has the independently finite original native
homology.  This constructs the frontier recursively, without taking
a frontier or its cohomology comparison as an external input.

At the requested degree, replace the flat frontier by its constructed
finite free middle term.  One further step supplies a finite free
incoming term.  The upper comparison shows that this last change
retains the requested native homology.  The outgoing term was already
finite free from the preceding stage.  Composing the actual native
homology isomorphisms gives the asserted arbitrary-scalar comparison.
\end{proof}

\begin{corollary}[Actual chain homology specializes on a principal open]
\label{cor:kn-derived-chain-homology-specialization}
\uses{thm:kn-derived-finite-free-triples,cor:kn-specialization-original-free-native-homology}
\lean{EnvelopeProperty.Vanishing.SpecializationDerivedFiniteApproximation.exists_chain_homology_zero_neighborhood,EnvelopeProperty.Vanishing.SpecializationDerivedFiniteApproximation.generic_chain_homology_isZero_of_closed_fibre}
\leanok
If the native degree-$q$ homology of the actual scalar-extended
complex vanishes at an original maximal ideal $\mathfrak p$, there
is $r\notin\mathfrak p$ such that it vanishes after every field
extension of coefficients in which $r$ is nonzero.  Over a domain,
it therefore vanishes over the original fraction field.
\end{corollary}
\begin{proof}
Construct the finite free triple, transfer the original residue
homology vanishing along its proved isomorphism, and apply finite
free homology specialization with the standard original bases.
Transfer the resulting zero groups back along the constructed
isomorphisms.  A nonzero element of the domain remains nonzero
in its actual fraction field.
\end{proof}

\begin{remark}[Relation to ordered section complexes]
\label{rmk:kn-derived-finite-free-triples-geometric-boundary}
The original-bound reversal and positive-degree natural-cochain
comparison are constructed in
Lemma~\ref{lem:kn-derived-cochain-reversal} and
Theorem~\ref{thm:kn-derived-natural-cochain-specialization}.
The actual geometric fibre cohomology interpretation is proved
separately in Lemma~\ref{lem:kn-derived-geometric-fibre-cohomology}.
\end{remark}
